% !TEX root = ../../../main.tex

\subsection{特征方案与模型输入设计}

句子视频识别任务在特征方案上没有直接沿用实时单词识别模型中的 30 帧 166 维输入，而是采用 20 帧增强特征作为词级分类模型的输入。
这样设计的主要原因在于，句子视频识别面对的是已经录制完成的短视频，视频中可能包含多个连续词汇、动作过渡帧和不同长度的词段。
如果仍然使用较长的 30 帧窗口，单个窗口更容易覆盖多个词汇或较长过渡区域，不利于后续进行词段划分。


\WhuAutoFigure{0.82\textwidth}{0.42\textheight}{figures/chapter06/fig6-8-plus-feature-232.png}{句子视频识别 $20\times232$ plus 特征构成示意图}{fig:chapter06-plus-feature-232}


\begin{longtable}{L{0.22\textwidth} L{0.12\textwidth} L{0.34\textwidth} L{0.24\textwidth}}
  \caption{句子视频识别 plus 特征维度构成表}
  \label{tab:chapter06-plus-feature-dim}\\
  \toprule
  特征部分 & 维度 & 组成 & 作用 \\
  \midrule
  base\_166 & 166 & 左手 78 + 右手 78 + 姿态 10 & 基础手型与上肢姿态 \\
  static\_plus\_28 & 28 & 姿态点、手中心、手框、手存在状态、距离 & 单帧空间增强 \\
  dynamic\_delta\_38 & 38 & 姿态点与手部代表点相邻帧差分 & 动作变化描述 \\
  合计 & 232 & 166 + 28 + 38 & 单帧 plus 特征 \\
  \bottomrule
\end{longtable}

其中，base\_166 是句子视频识别特征中的基础部分。
需要说明的是，该部分在维度组织上与实时单词识别中的 166 维输入相近，都是由左右手手型特征和上肢姿态辅助特征组成；
但二者的关键点提取链路并不完全相同。
句子视频识别部分直接使用 MediaPipe Holistic 全身模型对视频帧进行处理，并从同一次 Holistic 结果中读取 pose\_landmarks、left\_hand\_landmarks 和 right\_hand\_landmarks，再构造后续特征。

static\_plus\_28 是句子视频识别中额外增加的静态增强特征，用于补充单帧内更丰富的空间关系。
该部分包括鼻尖、左右肩、左右肘、左右腕等上半身关键点相对于肩中心的归一化坐标，两只手的中心位置，两只手的外接框宽高，左右手是否存在的标志，以及手腕到身体中心、手腕到手部中心的距离等信息。

dynamic\_delta\_38 是用于描述动作变化的动态差分特征。
系统先从每帧中提取一个 38 维动态源向量，该向量由 9 个上半身姿态点和左右手各 5 个代表性手部点组成，每个点使用二维坐标，共 19 个点、38 维。
随后，系统计算当前帧动态源向量与上一帧动态源向量之间的差值，形成动态差分特征。
